% Generated draft: requires mathematical and source review.
\documentclass[11pt]{article}
\usepackage[margin=25mm]{geometry}
\usepackage{fontspec}
\setmainfont{Latin Modern Roman}
\usepackage{amsmath,amssymb,mathtools,unicode-math}
\setmathfont{Latin Modern Math}
\usepackage{enumitem,needspace,xurl,hyperref}
\setlength{\parindent}{0pt}
\setlength{\parskip}{5pt}
\setlength{\emergencystretch}{3em}
\newcommand{\N}{\mathbb{N}}
\newcommand{\Z}{\mathbb{Z}}
\newcommand{\Q}{\mathbb{Q}}
\newcommand{\R}{\mathbb{R}}
\newcommand{\C}{\mathbb{C}}
\newcommand{\F}{\mathbb{F}}
\newcommand{\A}{\mathbb{A}}
\newcommand{\PP}{\mathbb{P}}
\newcommand{\E}{\mathbb{E}}
\newcommand{\eps}{\varepsilon}
\newcommand{\dd}{\,\mathrm{d}}
\begin{document}
\section*{UnsolvedMath ID: 700650}
\textbf{Source problem code:} catalogue entry
\subsection{Definitions and mathematical statement}
Let \(\overline{\Q}=\{z\in\C:z\text{ is algebraic over }\Q\}\), and let \(GL_n(K)\) denote the group of invertible \(n\times n\) matrices over a field \(K\). A group \(H\) is torsion-free if, for every \(h\in H\) and every integer \(m\geq 1\), the equality \(h^m=1\) implies \(h=1\). For a positive integer \(n\), let
\[
\mathcal{T}_n=\{H\leq GL_n(\overline{\Q}):H\text{ is torsion-free}\}.
\]
A group \(G\) is residually in \(\mathcal{T}_n\) if, for every \(g\in G\setminus\{1\}\), there exist \(H\in\mathcal{T}_n\) and a group homomorphism \(\varphi:G\to H\) such that \(\varphi(g)\neq 1\). The group \(H\) and the homomorphism may depend on \(g\).

The question attributed to G. A. Noskov is: for every positive integer \(n\), is every finitely generated torsion-free subgroup \(G\leq GL_n(\C)\) residually in \(\mathcal{T}_n\)? Equivalently, must every such \(G\) satisfy
\[
\forall g\in G\setminus\{1\}\quad
\exists\varphi:G\longrightarrow GL_n(\overline{\Q})
\quad\bigl(\varphi\text{ is a group homomorphism},\ 
\varphi(G)\text{ is torsion-free},\ 
\varphi(g)\neq 1\bigr)?
\]
The source leaves the quantification of \(n\) implicit; this draft uses the conventional reading that the question concerns every positive integer \(n\), with the same matrix dimension in the original group and the separating target groups. The malformed matrix-group delimiters in the supplied statement have been normalized.
\subsection{Short English statement}
For each positive matrix dimension \(n\), can every nonidentity element of a finitely generated torsion-free complex linear group of dimension \(n\) be detected by a homomorphism into a torsion-free group of matrices of the same dimension with algebraic entries? This is the conventional reading of the source's implicit dimension quantifier.
\subsection{Sources}
\begin{itemize}
\item Local UnsolvedMath registry, supplied record with permanent ID \(700650\), exact source problem code \texttt{catalogue entry}, title \emph{Kourovka Notebook Problem 10.35}; upstream revision \texttt{d371ee917ac9647fbb06cd6fc529deec6f3dda82}.
\item Supplied authoritative source attribution: \emph{The Kourovka Notebook (Unsolved Problems in Group Theory)}, original problem 10.35, 10th issue (1986), as recorded in the 21st-issue edition (2026), \texttt{arXiv:1401.0300v48}. The supplied collection description names E. I. Khukhro and V. D. Mazurov as editors.
\item Proposed by G. A. Noskov in 1986, according to the supplied record.
\item The registry labels the problem \emph{open}. This is supplied source metadata, not an independently verified assessment of its mathematical status. The cited edition was not independently consulted for this draft.
\end{itemize}
\end{document}
